\documentclass[12pt,a4paper,oneside]{article}
\usepackage[brazil]{babel}
\usepackage[utf8]{inputenc}
\usepackage{olymp}
\usepackage{color}
\usepackage[dvipsnames]{xcolor}
\usepackage{listings}

\setcounter{problem}{1}

\begin{document}

\begin{problem}{Movimentos da torre}{torre}{1}

Você começou a participar de um projeto que está desenvolvendo um tutorial de xadrez.
% O trecho de código abaixo imprime um tabuleiro $8 \times 8$.

% \lstset{language=C++,
%                 basicstyle=\ttfamily,
%                 keywordstyle=\color{Mulberry}\ttfamily,
%                 stringstyle=\color{Maroon}\ttfamily,
%                 commentstyle=\color{YellowGreen}\ttfamily,
%                 morecomment=[l][\color{magenta}]{\#},
%                 basewidth = {.48em}
% }
% \begin{lstlisting}
%     for(int i=1;i<=8;i++) {
%         for(int j=1;j<=8;j++)
%             cout << "-";
%         cout << endl;
%     }
% \end{lstlisting}



Sua primeira tarefa é mostrar os possíveis movimentos de uma torre. No xadrez, uma torre pode ser movida ortogonalmente no tabuleiro, ou seja, para casas que estejam na mesma linha ou na mesma coluna.

%\Notes
%
%Observações, se tiver.

\InputFile

A entrada contém dois números $L$ e $C$, que indicam a linha e a coluna da posição atual da torre. As linhas são numeradas de cima para baixo e as colunas da esquerda para direita. Restrições, $1 \leq L, C \leq 8$.


\OutputFile

Seu programa deve produzir uma saída como a mostrada abaixo, onde `\texttt{T}' indica a posição atual da torre, `\texttt{x}' (minúsculo) indica uma casa para onde ela pode se mover e `\texttt{-}' uma casa para onde ela não pode se mover.

\Example

\begin{example}
\exmp{
2 4
}{
---x----
xxxTxxxx
---x----
---x----
---x----
---x----
---x----
---x----
}%
\end{example}

\begin{example}
\exmp{
8 2
}{
-x------
-x------
-x------
-x------
-x------
-x------
-x------
xTxxxxxx
}%
\end{example}



% \begin{example}
% \exmp{
% 2 4
% }{
% \hspace*    1 2 3 4 5 6 7 8
% \hspace*.  ----------------
% 1| - - - x - - - -
% 2| x x x T x x x x
% 3| - - - x - - - -
% 4| - - - x - - - -
% 5| - - - x - - - -
% 6| - - - x - - - -
% 7| - - - x - - - -
% 8| - - - x - - - -
% }%
% \end{example}

% \begin{example}
% \exmp{
% 2 4
% }{
% \hspace*    1 2 3 4 5 6 7 8
% \hspace*.  ----------------
% 1| - - - x - - - -
% 2| x x x T x x x x
% 3| - - - x - - - -
% 4| - - - x - - - -
% 5| - - - x - - - -
% 6| - - - x - - - -
% 7| - - - x - - - -
% 8| - - - x - - - -
% }%
% \end{example}

%Comentários sobre o exemplo, se necessário.

\end{problem}

\end{document}
